%==============================================================================
% 2. INTRODUCCIÓN
%
% Pauta: "Deberán tomarse los siguientes puntos:
%   - Formulación del Marco Teórico, incluyendo antecedentes bibliográficos.
%   - Planteamiento del problema.
%   - Formulación de hipótesis.
%   - Objetivo(s) general(es) y específicos."
%==============================================================================

\chapter{Introducción}

\section{Marco teórico y antecedentes bibliográficos}

Desarrolla aquí el marco teórico con sus antecedentes bibliográficos. Recuerda
las dos formas de citar que exige la pauta: como parte de la frase,
\textcite{Dayton1984} describieron la dinámica de los bosques de
\especie{Macrocystis}; o entre paréntesis al final de la idea, la estabilidad
de estas comunidades depende de la frecuencia de perturbación
\parencite{Dayton1984}.

Cuando la referencia tiene dos autores se anotan ambos
\parencite{BeckerOrtiz1968}; cuando tiene más de dos, va el primer autor
seguido de \extranjero{et al.} \parencite{Becker1968}.

Los nombres científicos y las expresiones greco-latinas van en cursiva: con
\verb|\especie{...}| se obtiene \especie{Macrocystis pyrifera}. Lo mismo vale
para las palabras en idioma extranjero que no tengan equivalente en español:
\verb|\extranjero{...}| da \extranjero{upwelling}.

\section{Planteamiento del problema}

Expón aquí el vacío de conocimiento que tu tesis aborda y por qué importa.

\section{Hipótesis}

Enuncia la hipótesis de trabajo en una frase contrastable. Por ejemplo:

\begin{quote}
\singlespacing
La intensificación del viento favorable a la surgencia durante el período
2000--2024 se asocia a una disminución significativa de la temperatura
superficial del mar en la zona costera de Valdivia.
\end{quote}

\section{Objetivos}

\subsection{Objetivo general}

Evaluar \dots

\subsection{Objetivos específicos}

\begin{enumerate}[label=\arabic*., leftmargin=*, itemsep=0.3\baselineskip]
  \item Caracterizar \dots
  \item Cuantificar \dots
  \item Determinar \dots
\end{enumerate}
